\chapter[\emj{🧩}~Componentes Fuertemente Conexas (SCC · Kosaraju)]{\emj{🧩}~Componentes Fuertemente Conexas (SCC · Kosaraju)}

\begin{dsaquees}

Grupos de amigos en una red social 👥 donde "seguir" es de un solo
sentido. Una Componente Fuertemente Conexa es un grupo donde TODOS
pueden llegar a TODOS siguiendo flechas: desde cualquier persona del
grupo puedes alcanzar a cualquier otra y regresar.

Kosaraju encuentra estos grupos con un truco astuto: recorre el grafo,
luego INVIERTE todas las flechas y lo recorre de nuevo. Los nodos que
siguen "atrapados juntos" incluso con las flechas al revés forman una
componente fuertemente conexa.

\end{dsaquees}

\begin{dsaotraforma}

En un grafo dirigido, una Componente Fuertemente Conexa (Strongly
Connected Component) es un subconjunto máximo de nodos donde existe un
camino dirigido entre cualquier par de ellos en ambos sentidos.

El algoritmo de Kosaraju las halla en dos pasadas de DFS:
\begin{enumerate}
\item DFS sobre el grafo original, apilando cada nodo cuando TERMINA
      (orden de finalización, como en topo-sort).
\item Transponer el grafo (invertir todas las aristas) y hacer DFS
      sacando nodos de la pila; cada árbol de este segundo DFS es una SCC.
\end{enumerate}

Alternativa de una sola pasada: el algoritmo de Tarjan (usa índices de
descubrimiento y \verb|low-link|), más eficiente pero más difícil de
recordar bajo presión.

\end{dsaotraforma}

\begin{dsadiagrama}
\begin{verbatim}
Grafo dirigido:

  A ──> B ──> C ──> D
  ^     │     ^     │
  └─────┘     └─────┘

  A <-> B forman ciclo   -> SCC {A, B}
  C <-> D forman ciclo   -> SCC {C, D}

Paso 1 (DFS original): apila por orden de finalizacion
Paso 2 (grafo transpuesto): cada DFS-tree = una SCC

Resultado: {A, B}  y  {C, D}
Grafo condensado (cada SCC = 1 super-nodo) SIEMPRE es un DAG.
\end{verbatim}
\end{dsadiagrama}

\begin{dsacomofunciona}
\begin{verbatim}
LAS DOS PASADAS de Kosaraju

Grafo:  A<->B → C<->D

PASADA 1 (DFS en original, apilar al TERMINAR):
  dfs(A)→dfs(B)→dfs(C)→dfs(D)
  termina: D, luego C, luego B, luego A
  PILA (tope arriba):  [ A ]
                       [ B ]
                       [ C ]
                       [ D ]

TRANSPONER (invertir todas las flechas):  A<->B ← C<->D

PASADA 2 (sacar de la pila, DFS en transpuesto):
  saca A → dfs alcanza {A,B}  ........ SCC #1 = {A, B}
  saca B → ya visitado, salta
  saca C → dfs alcanza {C,D}  ........ SCC #2 = {C, D}
  saca D → ya visitado, salta

Resultado: 2 componentes fuertemente conexas.
La transposición "corta" los puentes entre SCC y aísla cada grupo.
\end{verbatim}
\end{dsacomofunciona}

\begin{lstlisting}[language=C++]
KOSARAJU (dos pasadas de DFS)
1. DFS en el grafo original; al terminar cada nodo, apílalo.
2. Construye el grafo transpuesto (invierte todas las aristas).
3. Saca nodos de la pila; por cada no visitado, un DFS en el
   transpuesto marca una SCC completa.

📝 PSEUDOCÓDIGO
──────────────────────────────────────────────────────────────

función kosaraju(V, adj):
    // Pasada 1: orden por finalizacion
    PARA cada u no visitado: dfs1(u) -> apila u al terminar
    adjT = transponer(adj)
    // Pasada 2: sacar de la pila sobre el transpuesto
    MIENTRAS pila no vacia:
        u = pila.pop()
        SI u no visitado_2:
            dfs2(u, adjT)      // marca una SCC entera
            sccCount++
    return sccCount

💻 CÓDIGO C++
──────────────────────────────────────────────────────────────

#include <vector>
#include <stack>
using namespace std;

void dfs1(int u, vector<vector<int>>& adj, vector<bool>& vis,
          stack<int>& st) {
    vis[u] = true;
    for (int v : adj[u]) if (!vis[v]) dfs1(v, adj, vis, st);
    st.push(u);                       // apila al TERMINAR
}

void dfs2(int u, vector<vector<int>>& adjT, vector<bool>& vis,
          vector<int>& comp) {
    vis[u] = true; comp.push_back(u);
    for (int v : adjT[u]) if (!vis[v]) dfs2(v, adjT, vis, comp);
}

int kosaraju(int V, vector<vector<int>>& adj) {
    vector<bool> vis(V, false);
    stack<int> st;
    for (int i = 0; i < V; i++) if (!vis[i]) dfs1(i, adj, vis, st);

    // Transponer
    vector<vector<int>> adjT(V);
    for (int u = 0; u < V; u++)
        for (int v : adj[u]) adjT[v].push_back(u);

    fill(vis.begin(), vis.end(), false);
    int scc = 0;
    while (!st.empty()) {
        int u = st.top(); st.pop();
        if (!vis[u]) {
            vector<int> comp;
            dfs2(u, adjT, vis, comp);  // una SCC
            scc++;
        }
    }
    return scc;
}

⏱️ COMPLEJIDAD (Big-O)
──────────────────────────────────────────────────────────────

  Algoritmo   │ Tiempo      │ Espacio
  ────────────┼─────────────┼──────────────
  Kosaraju    │ O(V + E)    │ O(V + E)
  Tarjan      │ O(V + E)    │ O(V)

* Tarjan hace una sola pasada; Kosaraju hace dos pero es más intuitivo.
\end{lstlisting}

\begin{dsaentrevistas}

✅ Grafo condensado = DAG: si colapsas cada SCC en un solo nodo, el
   grafo resultante nunca tiene ciclos. Sirve para aplicar topo-sort
   encima. Frase que impresiona.

💡 "Critical Connections" / puentes (LeetCode 1192): variantes de
   conectividad que se resuelven con low-link estilo Tarjan.

⚠️ SCC es solo para grafos DIRIGIDOS. En no dirigidos el concepto
   equivalente es "componentes conexas" simples (un DFS/BFS o Union-Find).

🔑 Recuerda el orden de Kosaraju: primero apilas por finalización en el
   grafo original, luego DFS en el TRANSPUESTO. Invertir ese orden lo
   rompe.

\end{dsaentrevistas}

\begin{dsaresumen}

• SCC: grupo máximo de nodos mutuamente alcanzables (grafo dirigido).
• Kosaraju: DFS original (apilar al terminar) + DFS en el transpuesto.
• Cada árbol del segundo DFS es una SCC.
• Tarjan: una sola pasada con low-link (más rápido, más difícil).
• Ambos O(V + E).
• Condensar las SCC produce siempre un DAG.
════════════════════════════════════════════════════════════════

\end{dsaresumen}
